\documentclass[10pt]{article}
\usepackage[margin=.84in]{geometry}
\usepackage{amsmath,amssymb,amsthm,booktabs,hyperref}
\hypersetup{colorlinks=true,urlcolor=blue,linkcolor=black,citecolor=black}
\newtheorem{theorem}{Theorem}
\newtheorem{lemma}[theorem]{Lemma}
\newtheorem{proposition}[theorem]{Proposition}
\newcommand{\F}{\mathbb F}
\title{Obstructions to cyclic polynomial labels over fields}
\author{Blake Stone's mathematics project\\\small Proof developed with OpenAI Codex}
\date{November 26, 2026}
\begin{document}\maketitle
\begin{abstract}
Zhi-Wei Sun asked whether every finite field can be circularly ordered so that the labels $P(a_i)+a_{i+1}$ exhaust the field for each polynomial of degree less than $q-1$ outside the exceptional family $c-X$. He also asked for distinct labels on arbitrary finite subsets. Both statements fail for $P=X^2$. The first fails on the field of eight elements, the second already on a three-point set over the rationals and in every characteristic greater than three. We give hand proofs and an exhaustive classification of all admissible polynomial functions on fields of order at most eight. The companion paper proves the original prime-field assertion for $X^2$ and classifies the moduli on which it holds over $\mathbb Z/n\mathbb Z$.
\end{abstract}

Sun's conjectures of 30 September 2013 appear in \cite{sun,oeis}. The first general assertion asks for a cycle on $\F_q$ under $\deg P<q-1$. The subset assertion concerns a field $F$ of order greater than seven, a set $A\subset F$ with $|A|>2$, and a polynomial with $\deg P<p-1$ in positive characteristic $p$. Both exclude $P(X)=c-X$. Write $\sigma(a_i)=a_{i+1}$ with cyclic indices. The label function is $P+\sigma$, where addition is pointwise.

\section{A quadratic obstruction over eight elements}
For an additive group $V$, a \emph{complete mapping} is a permutation $h$ of $V$ such that $x\mapsto x+h(x)$ is also a permutation. Bors and Wang \cite{bors} determined the permutation groups generated by complete mappings of finite fields, including the exceptional small orders. We give a direct parity argument in the two relevant dimensions.

\begin{lemma}\label{lem:even}
Every complete mapping of the additive group of $\F_4$ or $\F_8$ is even.
\end{lemma}
\begin{proof}
Regard the field as a vector space $V$ over $\F_2$. Translating the output by $h(0)$ preserves completeness, makes $h(0)=0$, and does not change sign. A nonidentity translation on $V$ consists of $|V|/2$ transpositions, an even number. Now $h$ has no nonzero fixed point, since $x+h(x)=0$ only for $x=0$. Nor can $h$ exchange two elements $x,y$, as $x+h(x)=x+y=y+h(y)$. On the three nonzero vectors of $\F_4$ this forces a three-cycle.

On the seven nonzero vectors of $\F_8$, the only alternatives are a seven-cycle and a three-cycle together with a four-cycle. Suppose the latter. The three edge labels in the three-cycle on $a,b,c$ are $a+b,b+c,c+a$ and are the three nonzero vectors of a plane $H$. Hence $a,b,c$ lie in the same coset of $H$. All four labels along the four-cycle must lie outside $H$, forcing its vertices to alternate between the two cosets. If $a,b,c\in H$, its four vertices all lie outside $H$. If $a,b,c\notin H$, its vertices consist of three in $H$ and one outside. Neither distribution permits alternation. Thus the seven nonzero vectors form a single cycle, of even sign.
\end{proof}

\begin{theorem}\label{th:reject}
Let $q\in\{4,8\}$ and let $P$ be any even permutation of $\F_q$. There is no circular ordering for which $P(a_i)+a_{i+1}$ are pairwise distinct. In particular $P=X^2$ fails on $\F_8$.
\end{theorem}
\begin{proof}
If such a cycle $\sigma$ existed, $h=\sigma\circ P^{-1}$ would be a permutation and
\[
 x+h(x)=(P+\sigma)(P^{-1}(x))
\]
would be a permutation. Lemma~\ref{lem:even} gives $\operatorname{sgn}(h)=1$, whereas a cycle of even length composed with an even permutation has negative sign. On $\F_8$ the squaring map is an automorphism of order three, so it is even. Its polynomial degree is $2<7$ and it is outside Sun's excluded family.
\end{proof}

\section{A three-point obstruction in characteristic zero and odd characteristic}
\begin{proposition}
Over any field $F$ of characteristic zero or $p\ge5$, the three elements $a=-5/4$, $b=1/4$ and $c=7/4$ are distinct. The labels for $P=X^2$ repeat in both circular orientations of this set.
\end{proposition}
\begin{proof}
The scaled differences are $4(b-a)=4(c-b)=6$ and $4(c-a)=12$. The two circles, up to rotation, are $(a,b,c)$ and $(a,c,b)$. The first has $a^2+b=b^2+c=29/16$ and the second has $a^2+c=c^2+b=53/16$. Thus both have a repeated label, in every stated characteristic. Taking $F=\mathbb Q$ gives the required field size, as do all finite fields of these characteristics having more than seven elements. In positive characteristic the degree satisfies $2<p-1$.
\end{proof}

\section{All polynomial functions through field order eight}
The following finite computation gives the converse of Theorem~\ref{th:reject} on these two fields, even allowing arbitrary admissible polynomials with repeated values.
\begin{theorem}
For $q=4,8$ and $\deg P<q-1$, a full cycle with distinct labels exists exactly when $P$ is not an even permutation of the field. For $q=3,5,7$, it exists exactly when $P$ is outside the family $c-X$. For $q=2$ every constant works.
\end{theorem}
\begin{proof}
Every field function has a unique polynomial representative of degree at most $q-1$. For $q>2$, summing a monomial over the field gives zero in degrees $0,\ldots,q-2$ and $-1$ in degree $q-1$. Thus the condition in the theorem is precisely $\sum_x P(x)=0$. Translating $P$ by a constant preserves distinctness of labels, so impose $P(0)=0$. The other $q-2$ freely chosen values determine the last one, giving $q^{q-2}$ normalized tables.

Enumerate every list beginning at zero and containing every field element once. Its successor map is one of the $(q-1)!$ distinct full cycles $\sigma$. For each such $\sigma$, enumerate all $(q-1)!$ permutations $\lambda$ satisfying $\lambda(0)=\sigma(0)$. Mark the normalized table $P(x)=\lambda(x)-\sigma(x)$. Every marked table works, with the cycle considered, and every solution appears when its cycle and its labels are visited. Every difference of these two permutations has sum zero. For the extensions of $\F_2$ the field additions and subtractions are just bitwise exclusive-or in a fixed additive basis; all other runs use prime-field arithmetic. The accompanying standard C++17 program performs the complete enumeration.
\begin{center}
\begin{tabular}{rrrrrr}\toprule
$q$&3&4&5&7&8\\\midrule
Normalized tables&3&16&125&16,807&262,144\\
Unmarked&1&3&1&1&2,520\\\bottomrule
\end{tabular}
\end{center}
At odd orders $P=-X$ is always unmarked, since a full cycle is fixed-point-free and $\sigma(x)-x$ then misses zero. At orders four and eight the even normalized permutations number $(q-1)!/2$, exactly the unmarked counts. Translating by any constant preserves permutation parity in those two fields. These facts identify every exception. The assertion at order two is immediate since translating the unique two-cycle preserves bijectivity.
\end{proof}

\begin{thebibliography}{9}
\bibitem{sun} Z.-W. Sun, \emph{Determinants, Permutations and Additive Combinatorics}, lecture slides, conjecture dated 30 September 2013. \url{https://www.mathnet.ru/PresentFiles/21090/sun.pdf}.
\bibitem{oeis} Z.-W. Sun, OEIS A229827, September 2013. \url{https://oeis.org/A229827}.
\bibitem{bors} A. Bors and Q. Wang, \emph{Compositions and parities of complete mappings and of orthomorphisms}, J. Combin. Theory Ser. A \textbf{196} (2023), 105723. \url{https://arxiv.org/abs/2207.09642}.
\end{thebibliography}
\end{document}
